\documentclass[11pt]{article}
\usepackage{amsmath,amssymb,amsthm,amsfonts,verbatim,fancyhdr,graphicx}
\usepackage{pdfpages}
\usepackage{enumerate}
\usepackage{microtype}
\usepackage{url}

\addtolength{\evensidemargin}{-.5in}
\addtolength{\oddsidemargin}{-.5in}
\addtolength{\textwidth}{0.8in}
\addtolength{\textheight}{0.6in}
\addtolength{\topmargin}{-.3in}

\theoremstyle{plain}
\theoremstyle{definition}
\newtheorem{theorem}{Theorem}[section]
\newtheorem{corollary}[theorem]{Corollary}
\newtheorem{proposition}[theorem]{Proposition}
\newtheorem{fact}[theorem]{Fact}
\newtheorem{lemma}[theorem]{Lemma}
\newtheorem{claim}[theorem]{Claim}
\newtheorem{conjecture}[theorem]{Conjecture}
\newtheorem{question}{Question}
\newtheorem{definition}[theorem]{Definition}
\newtheorem{example}[theorem]{Example}
\newtheorem{exercise}[theorem]{Exercise}
\newtheorem{challenge}[theorem]{Challenge Problem}

\newcommand{\Q}{\mathbb{Q}}
\newcommand{\R}{\mathbb{R}}
\newcommand{\Z}{\mathbb{Z}}
\newcommand{\C}{\mathbb{C}}

% \DeclareMathOperator{\sin}{sin}   %% just an example (it's already defined)
\DeclareMathOperator{\norm}{norm}
\renewcommand{\baselinestretch}{1.2}
\pagestyle{fancy}
\fancyhf{ }
\lhead{MAT 128B: Project 1}
\rhead{M. Fiacco C. Palomo}
\begin{document}
	

	\begin{titlepage}
	\begin{center}
			\vspace*{1cm}
			\textbf{MAT 128B Project 1: }\linebreak
			\textbf{Using Iteration Methods to Understand Fractal Geometry}\\
			\vspace{1.0cm}
			{Melissa Fiacco \& Carlos Palomo}\linebreak
			Date: 19 February 2018\linebreak	
		\end{center}
	\end{titlepage}

\subsection*{Introduction}
In this project we are going to implement a series of computer programs that use iteration methods to generate different figures on on the plane. We will be focusing on the Filled Julia set, Julia set, and the Mandelbrot Set: The Julia set is a set of complex numbers that do not converge to any limit when a given mapping is repeatedly applied to them, and the Mandelbrot Set is a particular set of complex numbers that has a highly convoluted fractal boundary when plotted representing convergence of the Julia set. Both sets produce interesting fractal geometry in the complex plane, which we will be able to generate using our iteration methods.

\subsection*{Part I: An introduction to Fractals}
$ \phi(z)=z^2$ can be transformed into a fixed point problem to predict its behavior based on the value z.If $ |z_0| \le 1$,the orbit remains bounded and is referred to as the Filled Julia set. We will not find a singular fixed point, but the Julia set will find a sequences points mapped closely to themselves (never repeating) infinitely in an orbit contained within unit circle. Note $ z_0 \le e^i = cos(1) + isin(1) $which describes the Filled unit circle in the complex plane. This process is what creates the fractal- a pattern occurs as the orbit is created.
Implementing the program on page 100, Julia.m, we produce the following figure, replicating figure 4.13 in the book:

\begin{center}
	\includegraphics*[scale = 0.5]{Plot1.png}
\end{center}

\subsection*{Part II: Generate (and plot) other examples changing the value of c in the function}

We can generate different fractal patterns when we plot $ \phi(z) = z^2 + c $ using various c values (see below for various examples). However, we must be careful to restrict $ |z_0| \le 2 $; If we do not make this restriction, the orbit is unbounded and $ z_0 $ is not contained in the Julia set. This is due to $ z_0 $ existing outside of the unit circle and$ \phi(z) = z^2 + c $ growing outside of the unit circle, spiraling out; The computer program will reject our value, and end the program.\\

If we change our $ z_0 $ values, we shift where the fractal is centered. It is possible to choose an initial value that will cause the program to fail, so we must be careful.
As an example, we took our code from Part I, and changed zo from $ -0.8 $ and $ -1.8 $ to $ -0.9 $ and $ -1.5 $, respectively,  and the resulting fractal is now slightly shifted up and to the left, which reflects chosen $ z_0 $:

\begin{center}
	\includegraphics*[scale = 0.5]{Plot2.png}\vspace{5cm}
\end{center}

Examples of various c values and their graphs: \\

\begin{tabular}{c c c}
	c=0.36+0.1i                                   &c=-0.123-0.745i                         &c= -0.40+0.6i\\
	\includegraphics*[scale = 0.3]{Plot3.png} 
	& \includegraphics*[scale = 0.3]{Plot4.png} 
	&\includegraphics*[scale = 0.3]{Plot5.png} 
\end{tabular}

\subsection*{Part III: Constructing the Julia Set}
In Part iI we constructed the Filled Julia set. In this part, we have constructed the Julia set, the boundary of the Filled Julia set for $ \Psi = \pm \sqrt{w-c} $.  The Inverse Iterative method utilizes multiple properties of complex numbers by splitting w and c into real and imaginary parts such that $ z=w-c $, and then finds $ \sqrt{z} $.\\

\begin{tabular}{c c c}
c=0.36+0.1i                                   &c=-0.123-0.745i                         &c= -0.40+0.6i\\
\includegraphics*[scale = 0.25]{Plot10.png} 
& \includegraphics*[scale = 0.25]{Plot11.png} 
&\includegraphics*[scale = 0.25]{Plot12.png} 
\end{tabular}


\subsection*{Part IV: Computing the Fractal Dimension}
The fractal dimension D is the complexity of a fractal: it measures the roughness of an object. The higher the fractal dimension, the more complex, and more rapidly the object changes as its scaled. For example, if we split a square into 4 squares, there will be N=4 squares similar to A, and each side is $ \frac{1}{2} $its original length, so $ D=\frac{4}{2}=2 $. With fractals, we cannot calculate this as easily, thus we must use algorithms to estimate the dimension.We created two algorithms based off the paper \underline{On calculation of fractal dimension of images } by Ajay Kumar Bisoi and Jibitesh Mishra. The first method is the Reticular Cell Counting method and the second is Differential Box-Counting method. Using these to algorithms we attempted to find the fractal dimension for plot found in Part I. The Reticular Cell Counting method found a fractal dimension of 2 and the Differential Box-Counting method found a fractal dimension of 1. Applying the algorithms to the filled unit circle we found a fractal dimension of 1

\subsection*{Part V: Connectivity of the Julia Set }
We created a program which computes orb(0) to determine connectivity of the Julia set. We say divergence occurs if $ |z| > 100 $, i.e. orb(0) is unbounded, thus the Julia set is not connected. After 1500 iterations, if $  |z |< 100 $, but our function does not converge to fixed point, we assume orb(0) is bounded, and Filled Julia set is connected. We tested our function with various c values we used in part II. Based on the figures created in part II, we expected $ c=0.36+0.1i $ and $ c=-0.123-0.745i $ to be connected and $ c=-0.4+0.6i $ not connected, which is what our program gave us.\\

\includegraphics*[scale = 0.7]{Result1.png}\\
\includegraphics*[scale = 0.7]{Result2.png}\\
\includegraphics*[scale = 0.7]{Result3.png}

\subsection*{Part VI: Coloring Divergent Sets }
We can extend part V to create a color coded figure that tells us how long it takes for an orbit to diverge given a c value. Our color map is darker the higher number of iterations it takes to diverge for a range of z values, given c. We tested our c values from part II, and received the following images. Note, since we diverge when $| z|>100 $, we extended the domain and range to include all values within $ |z|<100 $. \\

\begin{tabular}{c c c}
	c=0.36+0.1i                                  &c=-0.123-0.745i                         &c= -0.40+0.6i\\
	\includegraphics*[scale = 0.3]{Plot6.png} 
	& \includegraphics*[scale = 0.3]{Plot7.png} 
	&\includegraphics*[scale = 0.3]{Plot8.png} 
\end{tabular}

\subsection*{Part VII: Newton’s Method in the Complex Plane}

If we apply similar methods from the coloring orbits program, we can find the fixed points of a complex function using Newton’s Method. We used the same colormap, with black meaning divergence after no iterations or small amount of iterations, and white implying convergence. The fixed points exist on the boundaries of the two colors. We tested our code for the function $ z^n +1 $ for various n.\\

\begin{tabular}{c c c c}
	n=2                                &n=3                         &n=5	& n=15\\
	\includegraphics*[scale = 0.23]{Plot13.png} 
	& \includegraphics*[scale = 0.23]{Plot14.png} 
	&\includegraphics*[scale = 0.23]{Plot15.png}
	&\includegraphics*[scale = 0.23]{Plot16.png}
\end{tabular}

\subsection*{Part VIII: The Mandelbrot Set }
For $  \phi(z)=z^2+c $, the Mandelbrot set shows where the Julia set is connected for various c values (the black region), and how long it takes for $ \phi(z)=z^2+c $ with $ z_0=0 $ to diverge.The lighter the color (away from black), the less iterations required to determine divergence. As we get closer to the black region, the darker the color is, and the longer it takes to determine divergence.The black region is where the color map is equivalently 0, meaning no divergence occurs

\begin{center}
	\includegraphics*[scale = 0.5]{Plot9.png}
\end{center}

\subsection*{Reference}
Greenbaum, Anne; Chartier, Timothy P.. \underline{Numerical Methods: Design, Analysis, and}\linebreak \underline{Computer Implementation of Algorithms} (p. 98). Princeton University Press. Kindle Edition. 
\pagebreak
\includepdf[scale = 0.70, pagecommand = \subsection*{Commit History}]{CommitHistory.pdf}

\end{document}